\begin{proof}
\begin{enumerate}
\item \textbf{Assumptions and degenerate cases.}  
We assume throughout that \(X(\omega)\ge0\) for every \(\omega\in\Omega\) and that \(a>0\).  
If \(\mathbb{E}[X]=0\), then \(X=0\) \(\mathbb{P}\)-a.s., whence \(\mathbb{P}(X\ge a)=0\) for all \(a>0\) and the desired inequality holds with equality.  Hence we may suppose \(\mathbb{E}[X]>0\).

\item \textbf{Indicator of the event \(\{X\ge a\}\).}  
Define the indicator random variable
\[
\mathbf 1_{\{X\ge a\}}(\omega)=
\begin{cases}
1, & \text{if }X(\omega)\ge a,\\[2pt]
0, & \text{if }X(\omega)< a .
\end{cases}
\]
Clearly \(0\le \mathbf 1_{\{X\ge a\}}(\omega)\le1\) for all \(\omega\).

\item \textbf{Pointwise inequality.}  
For each \(\omega\in\Omega\) we have
\[
\mathbf 1_{\{X\ge a\}}(\omega)\le\frac{X(\omega)}{a}. \tag{1}
\]
Indeed, if \(X(\omega)\ge a\) then the left–hand side equals \(1\) and \(1\le X(\omega)/a\);  
if \(X(\omega)<a\) then the left–hand side is \(0\) and \(0\le X(\omega)/a\) because \(X(\omega)\ge0\).

\item \textbf{Monotonicity of expectation.}  
Since expectation preserves the order of non‑negative random variables, (1) yields
\[
\mathbb{E}\!\big[\mathbf 1_{\{X\ge a\}}\big]\le \mathbb{E}\!\Big[\frac{X}{a}\Big].
\]

\item \textbf{Linearity of expectation.}  
The factor \(1/a\) is constant, so
\[
\mathbb{E}\!\Big[\frac{X}{a}\Big]=\frac{1}{a}\,\mathbb{E}[X]. \tag{2}
\]

\item \textbf{Expectation of an indicator equals probability.}  
For any event \(A\in\mathcal{F}\), \(\mathbb{E}[\mathbf 1_{A}]=\mathbb{P}(A)\).  Taking \(A=\{X\ge a\}\) gives
\[
\mathbb{E}\!\big[\mathbf 1_{\{X\ge a\}}\big]=\mathbb{P}(X\ge a). \tag{3}
\]

\item \textbf{Conclusion.}  
Combining (2) and (3) with the inequality of step 4 we obtain
\[
\mathbb{P}(X\ge a)\le\frac{\mathbb{E}[X]}{a},
\]
which is Markov’s inequality.

\item \textbf{Equality case.}  
Equality throughout the chain occurs precisely when equality holds in the pointwise estimate (1) for \(\mathbb{P}\)-almost every \(\omega\).  Thus we must have
\[
\mathbf 1_{\{X\ge a\}}(\omega)=\frac{X(\omega)}{a}\quad\text{a.s.}
\]
which forces
\[
X(\omega)=0\ \text{whenever }X(\omega)<a,\qquad
X(\omega)=a\ \text{whenever }X(\omega)\ge a.
\]
Hence \(X\) takes only the two values \(0\) and \(a\) (the trivial case \(X\equiv0\) being included).  Conversely, if \(X\) assumes only these two values, then \(\mathbb{P}(X\ge a)=\mathbb{P}(X=a)\) and \(\mathbb{E}[X]=a\,\mathbb{P}(X=a)\), so the bound is attained.

\end{enumerate}
\end{proof}